\chapter{Study of Existing Solutions}

\section{Introduction}
This chapter looks at tools and platforms that solve parts of the same problem space as MapIT. The point is not to collect names for the sake of completeness, but to understand how existing products treat documentation, incident traceability, collaboration, and knowledge reuse. That comparison gives the report a concrete baseline for the design choices made later.

\section{Methodology for Comparing Existing Solutions}
The comparison uses criteria that matter for an IT documentation and knowledge management platform. These criteria include structured organization, search quality, incident traceability, collaboration, access control, auditability, deployment flexibility, and licensing constraints. Using the same criteria for every product makes it possible to compare tools that were originally designed for very different audiences.

\section{Review of Existing Platforms}
Platforms such as ServiceNow, Jira Service Management, Confluence, GLPI, BookStack, Notion, SharePoint, and IT Glue show that the market already offers mature solutions for ticketing, documentation, and knowledge sharing. The common pattern is that each tool is strong in one direction and less complete in the others: service desk control, collaborative writing, enterprise content management, or flexible note keeping. For MapIT, the relevant question is whether a tool keeps infrastructure documentation and incident knowledge close together without turning the workflow into an oversized enterprise suite.

\section{Comparative Analysis}
The comparison table makes the differences visible at a glance. Enterprise tools are usually stronger in governance and workflow control, while lightweight tools are easier to adopt but weaker in traceability and structured access control. The conclusion is that MapIT is valuable precisely because it stays focused: it aims for clarity, controlled scope, and academic feasibility instead of reproducing the full breadth of a large commercial platform.

\begin{table}[H]
\centering
\small
\setlength{\tabcolsep}{3pt}
\renewcommand{\arraystretch}{1.12}
\begin{tabular}{|>{\raggedright\arraybackslash}p{2.35cm}|>{\raggedright\arraybackslash}p{3.85cm}|>{\raggedright\arraybackslash}p{3.85cm}|>{\raggedright\arraybackslash}p{3.85cm}|}
\hline
\textbf{Solution} & \textbf{Main strengths} & \textbf{Main limits} & \textbf{Relevance for MapIT} \\ \hline
\textbf{ServiceNow} & Mature ITSM workflows, governance, and automation. & Heavy, costly, and broader than the project scope. & Benchmark for rigor and access control. \\ \hline
\textbf{GLPI} & Open-source service desk with asset management. & Less oriented to relationship maps and reusable knowledge. & Useful for compact IT support comparison. \\ \hline
\textbf{Confluence} / BookStack & Strong structured documentation and collaboration. & Keeps documentation apart from incidents and dependencies. & Good reference for knowledge-base organization. \\ \hline
\textbf{Notion} & Flexible and easy for note-taking. & Weak structure, traceability, and IT-specific workflow control. & Shows limits of generic knowledge tools. \\ \hline
\textbf{MapIT} & Integrates assets, incidents, problems, relationships, search, and access control. & Deliberately narrower than enterprise suites. & Target solution of the study. \\ \hline
\end{tabular}
\caption{Feature comparison of existing solutions}
\label{tab:feature-comparison}
\end{table}

\section{Limitations of the Existing Approaches}
The reviewed approaches reveal several recurring limits. Some tools are expensive or too heavy for a compact academic deployment, while others are flexible but too generic to enforce a meaningful relationship between incidents and technical knowledge. In many cases, documentation is still separated from operational history, so the lesson learned from a previous incident is not naturally carried forward into the next one.

\section{Positioning of MapIT}
MapIT is positioned as a focused web application that covers the essential needs of infrastructure documentation and incident knowledge management without the overhead of a large enterprise suite. It is not meant to replace the market leaders. Its role is to show that a coherent solution can be tailored to the workflow of an IT team that needs controlled access, a clear relational model, and a direct connection between incidents and reusable knowledge.

\section{Chapter Conclusion}
This chapter shows that existing solutions address the problem only partially or at a broader scale than what the project requires. The comparison supports the decision to build MapIT as a compact, role-aware, and documentation-centric platform. The next chapter introduces the theoretical concepts that define the project domain.
